
\section{Related work}
\label{sec:related}


\paragraph{Heuristic Samplers for dLLMs.} Sampling in diffusion LLMs (dLLMs) \citep{llada2025, ye2025dream} has recently attracted significant attention, with much of the work in this area proposing heuristic approaches to improve the decoding process in a training-free manner. Throughout this paper, we focus on Fast-dLLM \citep{fastdllm2025} as a representative heuristic method, as it popularized confidence-thresholded sampling in dLLMs and demonstrated its crucial role in enabling faster inference in dLLMs compared to autoregressive models. Many recently proposed heuristics can be viewed as either reinterpretations or extensions of such confidence thresholding strategies \citep{ben2025accelerated, wei2025accelerating, hong2025wide, li2025diffusion, yu2025dimple}. Other notable heuristic approaches explore incorporating spatial \citep{huang2025pc} or temporal information \citep{wang2025time}, alternative confidence measures \citep{kim2025train}, or more explicit modeling of token dependencies \citep{azangulov2025parallel}. In this work, we aim to complement ongoing research on sampling heuristics by investigating whether effective sampling strategies can be learned directly via reinforcement learning.
Note that beyond improving the efficiency of unmasking in dLLMs, heuristics have also been proposed for remasking \citep{hong2025wide, dong2025saber} and for dynamically adjusting the generation length \citep{li2025beyond}, which we do not target in this work.

\paragraph{Reinforcement Learning Post-Training for dLLMs.} Early work on post-training diffusion LLMs via reinforcement learning includes d1 \citep{zhao2025d1}, which introduces a variant of GRPO tailored to dLLMs, and DiffuCoder \citep{gong2025diffucoder}, which focuses on enhancing the coding abilities of LLaDA-style models using RL. Most recent extensions aim to improve the quality of policy gradient estimators \citep{tang2025wd1, wang2025spg, lin2025boundary, rojas2025improving, zhu2025enhancing, wang2025revolutionizing, zhan2025principled}, and have demonstrated promising results in further enhancing the reasoning capabilities of dLLMs. The key distinction from our approach is that these methods use a fixed sampling strategy (e.g., high-confidence sampling), and the policy corresponds to the dLLM itself (or a LoRA-augmented version for efficiency). Closest to our work are DCOLT \citep{huang2025reinforcing}, which trains a separate unmasking module in addition to updating the base model via RL, and DiFFPO \citep{zhao2025diffpo}, which jointly learns an unmasking confidence threshold while updating the base model through RL. 
Differently to both these works, our primary goal is not in improving the reasoning abilities of the underlying dLLM; hence we keep the underlying dLLM fixed and focus solely on training a standalone policy with the aim of learning fast sampling while preserving the performance of the base model.
Finally, in concurrent work, Seed Diffusion \citep{song2025seed} proposes computation-aware RL as one of the key ingredients for achieving competitive efficiency among closed-source coding dLLMs.

\paragraph{Orthogonal efforts to accelerate dLLMs.} Besides the aforementioned heuristic-based sampling innovations, other efforts to improve inference efficiency in dLLMs \citep{ma2025dinfer} include KV caching \citep{jiang2025d, fastdllm2025}, (variants of) speculative decoding \citep{israel2025accelerating, campbell2025self, guo2025reviving}, training separate decoder modules during pretraining \citep{liu2024think, arriola2025encoder}, and diffusion forcing \citep{wang2025diffusion}, among others. Concurrent work explores distilling faster generation patterns directly into the base model \citep{chen2025dparallel}, or training a (small) separate unmasking network on top of the pretrained dLLM \citep{bansal2025enabling, bao2025learning}. Crucially, unlike our approach, where this training is done via RL, these works train unmasking modules by distilling generation trajectories from the base model. As a result, we expect that such approaches may still inherit limitations observed in dLLM sampling, such as difficulties in the non semi-AR regime (cf. \Cref{sec:back_heuristics}).

\paragraph{RL for Adaptive Compute} Since our sampling policies result in a variable number of sampling steps ($T - \hat{T}$) per input, our work connects to the broader literature on adaptive computation \citep{graves2016adaptive, teerapittayanon2016branchynet}. Notable examples of using RL to learn input-dependent compute policies include conditional computation with stochastic gating \citep{bengio2015conditional}, dynamic block skipping in residual networks \citep{wang2018skipnet}, and RL-based early-exit decisions for long chain-of-thought reasoning \citep{dai2025s}. To the best of our knowledge, our work is the first to explore learning adaptive policies using RL for the task of sampling in dLLMs. The closest related concurrent work is DiFFPO \citep{zhao2025diffpo}; however, unlike our approach--where the policy is learned end-to-end--DiFFPO learns to predict adaptive thresholds, which are then used in the same fashion as the fixed thresholds employed by Fast-dLLM \citep{fastdllm2025}.


